\documentclass{article}

\usepackage{hyperref}
\usepackage{fancyvrb}

\begin{document}

\noindent
\textbf{\Large Koji Sunami}
\hfill
\qquad
\href{https://kojisunami.github.io/index.html}{\textbf{Top}}
\qquad
\href{https://kojisunami.github.io/misc.html}{\textbf{Misc}}
\qquad
\href{https://kojisunami.github.io/cv-koji.pdf}{\textbf{CV}}



\noindent\rule{\textwidth}{0.4pt}


In this page, we’ll introduce many other non-math general stuffs. Some of them are my essays.

\noindent\rule{\textwidth}{0.4pt}
 

\textbf{Web Development With Latex}

I use Latex and github to design this website, and I’ll talk about the construction howto of the environment. 
The Latex compiler I’m using is Latexmlc, that outputs .html files from .tex file, 
while its CLI automation of the uploading files is done by github. The whole automation procedure is programmed on its Makefile given below:

\begin{verbatim}
HTML_PROJECTS = index random-thoughts misc
PDF_PROJECTS = cv-koji

.PHONY: all src

all: $(HTML_PROJECTS:%=%.html) $(PDF_PROJECTS:%=%.pdf) src

%.html : %.tex
    latexmlc $< --dest=$@ --css=style.css

%.pdf: %.tex
    pdflatex $*

src:
    make -C src

clean:
    rm *aux *log *out
    make -C src clean

run:
    git add .
    git status
    git commit -a -m "uploading"
    git push origin master
\end{verbatim}
 
Continuing from the main html web page, I separated some exposition pdf files on src/ directory where there is another Makefile linked
as follows:

\begin{verbatim}
PROJECTS = ainfty-geometry dgla-stack dt-moduli  what-math
.PHONY: all
all: $(PROJECTS:%=%.pdf)

%.pdf : %.tex
    pdflatex $*
    biber $*
    pdflatex $*
    pdflatex $*

clean :
    rm *aux *bbl *bcf *blg *log *out *xml
\end{verbatim}




\noindent\rule{\textwidth}{0.4pt}


\textbf{Algorithm of Finding University That Fits You}

Searching new environment that fits us isn’t easy, but making database of professors and their research interests 
will help our future career selection. There’s already mathjobs.org for postdoc jobs, 
but there seems to be no equivalent for the PhD admission, and I use python code with OpenAlex.
OpenAlex is a web database of academic network, where we can search the name of professors and their profile,
including their research papers and their coauthors, and we can use it from python library called pyalex from an automated way.

The procedure of finding the list of professors relevant to my interest is to find coauthor of my current advisor Yan Soibelman, 
so I’ll search the profile of my advisor by using pyalex library to extract all of his coauthors.
In addition, once we find his coauthors, we’ll find coauthors of coauthors. This is Erdos number 2.

\begin{verbatim}
from pyalex import Authors, Works
import csv
import os

#name = input(": ")
name = "Yan Soibelman"

def find_coauthor (name):
    author = Authors().search(name).get()[0]

     coauthors = set()

    for work in Works().filter(
        authorships={"author": {"id": author["id"]}}
    ).get():
        for a in work["authorships"]:
            if a["author"]["id"] != author["id"]:
                coauthors.add(a["author"]["display_name"])

    print(f"\n{author[’display_name’]} Coauthor:")

    #for name in sorted(coauthors):
    #    print(name)
    return sorted(coauthors)


# find and make list of professors
names = []
for name in find_coauthor(name):
    #print (find_coauthor(name))
    names.extend(find_coauthor(name))


# eliminate double counting from the list
names = list (dict.fromkeys(names))


# export data in csv
filename = "coauthors.csv"
i = 1
while os.path.exists(filename):
    filename = f"coauthors{i}.csv"
    i += 1

with open(filename, "w", newline="", encoding="utf-8") as f:
    writer = csv.writer(f)
    for name in names:
        writer.writerow([name])
\end{verbatim}


Now once we have created the list of professors, we’ll type each of their names with pyalex and gain a profile 
about the name of the university they belong, research interests, and the title of top 3 most cited research papers respectively. 
Once we finish collecting the data, we’ll output in csv. Here is another python code beginning from the professors’s name list:

\begin{verbatim}
from pyalex import Authors, Works
import csv
import os
## search the information of professors
def info_professor(name) :
    author = Authors().search(name).get()[0]
    institutions = author["last_known_institutions"]
    university = institutions[0]["display_name"] if institutions else "NONE"

    topics = author["topics"][:3]
    topic_names = [t["display_name"] for t in topics]

    # 
    works = Works().filter( author={"id": author["id"]}).sort(cited_by_count="desc").get()
    top_papers = works[:3]

    paper_names = [
        f"{w[’title’]} ({w[’cited_by_count’]} citations)"
        for w in top_papers
    ]

    return [author, institutions, university, " ".join(topic_names), "\n".join(paper_names)]

# Run functions
list_profile = []
for name in names:
    list_profile.append(info_professors(name))


# export data in csv
filename = "coauthors_profile.csv"
i = 1
while os.path.exists(filename):
    filename = f"coauthors_profile{i}.csv"
    i += 1

with open(filename, "w", newline="", encoding="utf-8") as f:
    writer = csv.writer(f)
    for profile in list_profile:
        writer.writerow([profile])
\end{verbatim}


Note finally in my environment (Debian WSL on Windows 11), the following commands seemed to be needed to run the python code:

\begin{verbatim}
sudo apt install python3-venv
python3 -m venv .venv
source .venv/bin/activate
pip install pyalex
\end{verbatim}



\noindent\rule{\textwidth}{0.4pt}


\textbf{How to Find Jobs}

First of all, I realized the existence of GTA job in US institute as PhD, 
because my professor recommended reading ”A Mathematician’s Survial Guide” 
\href{https://www.amazon.com/Mathematicians-Survival-Guide-Graduate-Development/dp/082183455X}
{https://www.amazon.com/Mathematicians-Survival-Guide-Graduate-Development/dp/082183455X}
during my undergraduate. 
This book tells us not only the academic system involving qualifying exams and thesis defense, 
but including PhD admission and postgraduation career path. I became interested in becoming PhD student partly because of this.
There is also a free online link "Postdoc Career Advice" 
\href{https://math.ou.edu/ kmartin/career-ad.html}
{https://math.ou.edu/ kmartin/career-ad.html}.

On the other hand,
in case we will drop out from PhD and in case we will not be able to find a job that is related to our specialization, 
there’s still tons of moderately paid good jobs like fisherman for hand-to-mouth.
Data Annotation https://dataannotation.tech (there are many other good sites) introduces part-time employment on online platform, 
where its duty is typically judging and fixing the error of AI decision by human critical thinking. 
Data Annotation introduces many data science but in particular, it introduces job for math proof judging, or if you are bilingual, 
I’d recommend translation of technical terms on online platform as a part-time job. 
There’s tons of stuffs that AI cannot accurately translate, so pay is good.

 
\noindent\rule{\textwidth}{0.4pt}

\textbf{Philosophy Of Research}

The goal of my research is not only to prove something but to make it user friendly, so we can actively think about it. 
If we encounter to the theorem, the theorem should give us new inspiration to make something new like a new conjecture 
or cultivating curiosity to another, and it’s not only its easiness of interpretation. 
I’d like to avoid ”so what” argument all the cost for people of any background, 
and my project is typically goal oriented when doing pure math, and I don’t like titling ”introduction to something”, 
because what we gain is knowledge not the speculation. 
As Sophie’s World says ”The only way we can find meaning in life is through philosophizing.”

why category why algebraic geomettry what geometry why algebra?

\end{document}
 


